% GENERATED FILE. Edit canonical lessons, then run npm run generate:books.
% canonical-source-hash: fnv1a64:c5526e01219e9a83
% canonical-lessons: FA-C351-sepah, FA-C351-nabard, FA-C351-piruzi, FA-C351-shekast-dadan, FA-C351-fath-kardan

\chapter{War and victory}
\label{ch:fa-war-and-victory}
\hlchaptermodality{351}

\begin{chapteropening}
\textbf{By the end of this chapter:} \emph{I can talk about an army, a battle, a victory, defeating an enemy and capturing a city.}

Complete the last lesson of chapter 351: I can talk about an army, a battle, a victory, defeating an enemy and capturing a city.
\end{chapteropening}

\section[\texorpdfstring{\fa{سپاه} (sepâh)}{سپاه (sepâh)}]{\fa{سپاه} (sepâh) — an army}

\label{lesson:FA-C351-sepah}



\begin{quote}
Before the new one: say the Persian for a castle, then the Persian for a crown.
\end{quote}

\subsection*{You'll want to know: \fa{سپاه}}
\textbf{\fa{سپاه}} (\emph{sepâh}) — \textquotedblleft{}an army\textquotedblright{}.

\begin{quote}
\fa{سپاه} \fa{شب} \fa{به} \fa{شهر} \fa{رسید}.
\end{quote}

\begin{quote}
\emph{Sepâh shab be shahr resid.}
\end{quote}

\begin{quote}
The army reached the city at night.
\end{quote}

Literary and military; \emph{lashkar} is a synonym.

\begin{cousinweb}
\emph{sepâh} is a native Persian word that EtymDB glosses \textquotedblleft{}army\textquotedblright{} without a parent; it is usually traced to Middle Persian \emph{spāh}.
\end{cousinweb}

\begin{grammarlens}[title={What you've built}]
One more word for telling, explaining and describing at B1.
\end{grammarlens}

\subsection*{Your turn}
\begin{itemize}
\raggedright
  \item \emph{Say it:} \emph{sepâh}
  \item \emph{Say it:} \emph{sepâh}, once more
  \item \emph{Say it:} a sentence of your own with \emph{sepâh}
  \item \emph{Recall:} say the Persian for a castle, then the Persian for a crown, then say \emph{sepâh} again
\end{itemize}

\begin{tcolorbox}[breakable,colback=teal!4,colframe=teal!35!black,title={Before you move on}]
What does \fa{سپاه} mean? (\textquotedblleft{}An army\textquotedblright{}.) Say it once more.
\end{tcolorbox}
\section[\texorpdfstring{\fa{نبرد} (nabard)}{نبرد (nabard)}]{\fa{نبرد} (nabard) — a battle, combat}

\label{lesson:FA-C351-nabard}



\begin{quote}
Before the new one: say the Persian for a crown, then the Persian for an army.
\end{quote}

\subsection*{You'll want to know: \fa{نبرد}}
\textbf{\fa{نبرد}} (\emph{nabard}) — \textquotedblleft{}a battle, combat\textquotedblright{}.

\begin{quote}
\fa{نبرد} \fa{سه} \fa{روز} \fa{طول} \fa{کشید}.
\end{quote}

\begin{quote}
\emph{Nabard se ruz tul keshid.}
\end{quote}

\begin{quote}
The battle went on for three days.
\end{quote}

Literary; in speech \emph{jang}.

\begin{cousinweb}
\emph{nabard} is a native Persian word that EtymDB glosses \textquotedblleft{}combat, battle\textquotedblright{} without a parent.
\end{cousinweb}

\begin{grammarlens}[title={What you've built}]
One more word for telling, explaining and describing at B1.
\end{grammarlens}

\subsection*{Your turn}
\begin{itemize}
\raggedright
  \item \emph{Say it:} \emph{nabard}
  \item \emph{Say it:} \emph{nabard}, once more
  \item \emph{Say it:} a sentence of your own with \emph{nabard}
  \item \emph{Recall:} say the Persian for a crown, then the Persian for an army, then say \emph{nabard} again
\end{itemize}

\begin{tcolorbox}[breakable,colback=teal!4,colframe=teal!35!black,title={Before you move on}]
What does \fa{نبرد} mean? (\textquotedblleft{}A battle, combat\textquotedblright{}.) Say it once more.
\end{tcolorbox}
\section[\texorpdfstring{\fa{پیروزی} (piruzi)}{پیروزی (piruzi)}]{\fa{پیروزی} (piruzi) — a victory}

\label{lesson:FA-C351-piruzi}



\begin{quote}
Before the new one: say the Persian for an army, then the Persian for a battle.
\end{quote}

\subsection*{You'll want to know: \fa{پیروزی}}
\textbf{\fa{پیروزی}} (\emph{piruzi}) — \textquotedblleft{}a victory\textquotedblright{}.

\begin{quote}
\fa{مردم} \fa{پیروزی} \fa{را} \fa{جشن} \fa{گرفتند}.
\end{quote}

\begin{quote}
\emph{Mardom piruzi râ jashn gereftand.}
\end{quote}

\begin{quote}
The people celebrated the victory.
\end{quote}

From \emph{piruz} \textquotedblleft{}victorious\textquotedblright{}. Also a win in sport.

\begin{cousinweb}
\emph{piruzi} is inherited from Middle Persian \emph{pērōzīh}, \textquotedblleft{}victory\textquotedblright{} (EtymDB).
\end{cousinweb}

\begin{grammarlens}[title={What you've built}]
One more word for telling, explaining and describing at B1.
\end{grammarlens}

\subsection*{Your turn}
\begin{itemize}
\raggedright
  \item \emph{Say it:} \emph{piruzi}
  \item \emph{Say it:} \emph{piruzi}, once more
  \item \emph{Say it:} a sentence of your own with \emph{piruzi}
  \item \emph{Recall:} say the Persian for an army, then the Persian for a battle, then say \emph{piruzi} again
\end{itemize}

\begin{tcolorbox}[breakable,colback=teal!4,colframe=teal!35!black,title={Before you move on}]
What does \fa{پیروزی} mean? (\textquotedblleft{}A victory\textquotedblright{}.) Say it once more.
\end{tcolorbox}
\section[\texorpdfstring{\fa{شکست} \fa{دادن} (shekast dâdan)}{شکست دادن (shekast dâdan)}]{\fa{شکست} \fa{دادن} (shekast dâdan) — to defeat, to beat}

\label{lesson:FA-C351-shekast-dadan}



\begin{quote}
Before the new one: say the Persian for a battle, then the Persian for a victory.
\end{quote}

\subsection*{You'll want to know: \fa{شکست} \fa{دادن}}
\textbf{\fa{شکست} \fa{دادن}} (\emph{shekast dâdan}) — \textquotedblleft{}to defeat, to beat\textquotedblright{}.

\begin{quote}
\fa{سپاه} \fa{کوچک} \fa{دشمن} \fa{را} \fa{شکست} \fa{داد}.
\end{quote}

\begin{quote}
\emph{Sepâh-e kuchak doshman râ shekast dâd.}
\end{quote}

\begin{quote}
The small army defeated the enemy.
\end{quote}

\emph{shekast} \textquotedblleft{}a break, a defeat\textquotedblright{} + \emph{dâdan} \textquotedblleft{}to give\textquotedblright{}. Present \emph{shekast midahad}.

\begin{cousinweb}
\emph{shekast} \textquotedblleft{}a break, a defeat\textquotedblright{} is the past stem of the native verb \emph{shekastan} \textquotedblleft{}to break\textquotedblright{}, and \emph{dâdan} \textquotedblleft{}to give\textquotedblright{} is inherited Persian: to give someone a defeat.
\end{cousinweb}

\begin{grammarlens}[title={What you've built}]
One more word for telling, explaining and describing at B1.
\end{grammarlens}

\subsection*{Your turn}
\begin{itemize}
\raggedright
  \item \emph{Say it:} \emph{shekast dâdan}
  \item \emph{Say it:} \emph{shekast dâdan}, once more
  \item \emph{Say it:} a sentence of your own with \emph{shekast dâdan}
  \item \emph{Recall:} say the Persian for a battle, then the Persian for a victory, then say \emph{shekast dâdan} again
\end{itemize}

\begin{tcolorbox}[breakable,colback=teal!4,colframe=teal!35!black,title={Before you move on}]
What does \fa{شکست} \fa{دادن} mean? (\textquotedblleft{}To defeat, to beat\textquotedblright{}.) Say it once more.
\end{tcolorbox}
\section[\texorpdfstring{\fa{فتح} \fa{کردن} (fath kardan)}{فتح کردن (fath kardan)}]{\fa{فتح} \fa{کردن} (fath kardan) — to conquer, to capture}

\label{lesson:FA-C351-fath-kardan}



\begin{quote}
Before the new one: say the Persian for a victory, then the Persian for to defeat.
\end{quote}

\subsection*{You'll want to know: \fa{فتح} \fa{کردن}}
\textbf{\fa{فتح} \fa{کردن}} (\emph{fath kardan}) — \textquotedblleft{}to conquer, to capture\textquotedblright{}.

\begin{quote}
\fa{سپاه} \fa{شهر} \fa{را} \fa{بعد} \fa{از} \fa{یک} \fa{ماه} \fa{فتح} \fa{کرد}.
\end{quote}

\begin{quote}
\emph{Sepâh shahr râ ba'd az yek mâh fath kard.}
\end{quote}

\begin{quote}
The army captured the city after a month.
\end{quote}

Also of a summit: \emph{qolle râ fath kard} \textquotedblleft{}he conquered the peak\textquotedblright{}. Present \emph{fath mikonad}.

\begin{cousinweb}
\emph{fath} is borrowed from Arabic \emph{fatḥ} (EtymDB), usually glossed \textquotedblleft{}an opening, a conquest\textquotedblright{}. \emph{kardan} \textquotedblleft{}to do\textquotedblright{} is inherited Persian.
\end{cousinweb}

\begin{grammarlens}[title={What you've built}]
One more word for telling, explaining and describing at B1.
\end{grammarlens}

\subsection*{Your turn}
\begin{itemize}
\raggedright
  \item \emph{Say it:} \emph{fath kardan}
  \item \emph{Say it:} \emph{fath kardan}, once more
  \item \emph{Say it:} a sentence of your own with \emph{fath kardan}
  \item \emph{Recall:} say the Persian for a victory, then the Persian for to defeat, then say \emph{fath kardan} again
\end{itemize}

\begin{tcolorbox}[breakable,colback=teal!4,colframe=teal!35!black,title={Before you move on}]
What does \fa{فتح} \fa{کردن} mean? (\textquotedblleft{}To conquer, to capture\textquotedblright{}.) Say it once more.
\end{tcolorbox}
